\subsubsection{Entanglement and Quantum Correlation}\label{sec:entanglement-and-quantum-correlation}

The EPR paradox highlights a fundamental feature of entangled systems: measurement outcomes may appear completely random to individual observers, yet exhibit strong correlations when compared afterward.

\highspace
In simple terms, Alice and Bob each observe \textbf{random outcomes} when measuring their respective qubits. However, when they later compare their results, they discover that these outcomes are \textbf{correlated in a way that appears difficult to reconcile with classical intuition}.

\highspace
\textcolor{Green3}{\faIcon{question-circle} \textbf{Alice and Bob obtain correlated results. Why is that surprising?}} At first glance, it might seem that correlated results are not surprising at all. After all, \hl{classical systems can also exhibit correlations}.

\highspace
\textcolor{Green3}{\faIcon{dice} \textbf{Classical example.}}
Suppose two coins are prepared in such a way that they always show the same face. If one observer sees Heads, they immediately know that the other coin must also be Heads. Likewise, if one observer sees Tails, the other must see Tails. This is an example of a \textbf{classical correlation}.

\highspace
In \hl{classical physics}, such \hl{correlations} are explained by assuming that the \hl{system possessed predetermined properties} before any observation was made. In our example, the two coins were already both Heads or both Tails before anyone looked at them.

\highspace
\textcolor{Green3}{\faIcon[regular]{lightbulb} \textbf{Einstein's intuition.}} Einstein believed that the \hl{correlations observed in entangled particles might be explained in exactly the same way}. According to this view, the particles could carry some form of \textbf{hidden information that predetermines the outcome of future measurements}. If such hidden information existed, the apparent mystery of entanglement would disappear: measuring one particle would simply reveal information that was already present from the beginning.

\highspace
Quantum entanglement exhibits a superficially similar phenomenon. Individual measurement outcomes appear random, yet they become strongly correlated when compared afterward.

\highspace
\textcolor{Red2}{\faIcon{exclamation-triangle} \textbf{The difference.}} Unlike classical correlations, which are completely determined by pre-existing properties of the system, \textbf{quantum correlations depend on the measurement basis chosen by the observers}. This means that \hl{changing how Alice and Bob perform their measurements can change the observed correlations in a way that has no classical analogue.}

\highspace
Therefore, the existence of correlations is not itself surprising. The \hl{real challenge is determining whether quantum correlations can be explained using the same classical intuition} that explains correlated coins, cards, or dice, or whether they \hl{represent an entirely new physical phenomenon.}

\highspace
This question leads directly to the distinction between \textbf{classical correlation} and \textbf{quantum correlation}.